\chapter{Conclusion}%
\label{ch:conclusion}

This thesis has involved modelling, parameter estimation, and control strategies for a tendon-driven hand exoskeleton intended to control both fingertip position and interaction forces. The system is intended for use in spinal-cord-injured patients, with the goal of improving hand function and independence. The main challenges addressed include the unknown and user-dependent dynamics of the fingers, the nonlinearities and friction in the bowden tube cable transmission, and the limited measurements at the hand. The work has focused on the design and simulation of control strategies, with the goal of providing a framework for future implementation and testing on the physical system. This includes modeling of a brushed DC motor, a planar finger, the cable transmission between the motors and the finger, as well as the design of a discrete-time current controller, a torque control scheme, and a model reference adaptive controller for position control.

A model of the brushed DC motor was obtained, including the electrical dynamics of the armature circuit and the mechanical dynamics of the rotor. Estimates for the motor parameters were calculated from the motor datasheet, and a method for estimating the remaining unknown parameters was developed and simulated. Further, a simplified planar model of the finger was derived using the Euler-Lagrange equations, including the effects of joint stiffness and damping. The parameters of the finger model were estimated empirically through simulations, to provide a natural looking motion of the finger. The cable connecting the motor to the finger was assumed to be sufficiently stiff, and modeled as a constraint relating the motor angle to the finger joint angles. 

A discrete-time current controller was designed for the brushed DC motor, using proportional feedback with a feedforward term to compensate for the \mbox{back-EMF}. Simulations showed that the controller was able to track the reference current very rapidly, compared with the slow mechanical dynamics of the motor. This supports the goal of using the armature current as a control input for both torque and position control. The stability of the discrete-time current controller was analyzed, and resulted in an upper bound for the proportional gain ensuring stability, based on the controller sampling time and some of the motor parameters. Simulations confirmed that the controller remained stable for gains below the upper bound. They also showed that a short sampling time is important for achieving good current tracking, where a longer sampling time led to steady-state errors and a slower response. Given the simplicity of the controller, simulations also showed it being robust to measurement noise.

Torque and force control were then considered using the current controller as its basis. Due to the generated motor torque being proportional to the current, controlling the current provided an indirect way to control the torque output of the motor. Considering the friction acting on the cable between the motor and the finger, the applied force at the fingertip won't be equal that which the motor pulls on the cable. Hence, integrated force sensors in the fingertips were considered. Using a proportional and integral (PI) feedback, the fingertip force sensors were shown through simulations to be able to combat the effects of friction and improve the tracking of the desired fingertip force, and should hence be considered for the real system. The PI-term could also be used to estimate the friction in the system, which could prove useful for monitoring the system, and potentially for diagnosing issues, such as wear and tear of the cables or bowden tubes.

For position control, given that the system is meant to be used by several different subjects, with different finger stiffness, retuning of the controller for each user would likely be required to provide a good user experience. To address this and to hopefully make the tuning process easier, a model reference adaptive control scheme was considered. Since the position of the fingertip is not directly measured, the motor angle was used as a proxy for the finger position, based on the assumption of a sufficiently stiff cable connection between the motor and the finger. Simulations showed that the MRAC scheme was able to make the motor angle track a reference model well over time. Initially it showed oscillatory behavior about the reference model. But over time, as the adaptive gains updated, these oscillations disappeared, and the controller was able to track the reference model very closely. However, the adaptive gains were also observed to continue growing slowly over time, and the error between the motor angle and the reference model never converged exactly to zero. This is most likely due to nonlinearities from the cable constraint, saturation in the armature reference current, and model structures that violate the ideal assumptions of the MRAC derivation. In practice, this signals that the MRAC scheme is rather to be used as an automatic tuning method for finding good fixed gains. That way the controller can freeze the adaptation gains after the error stops decreasing, or after a predefined time period. This was tested in simulations, and showed that the controller was able to track the reference model well, even with static gains found by the MRAC scheme. 

Considering that the controllers developed doesn't rely on any force measurements prior to the bowden tube, the force sensors currently mounted on the system can be considered redundant, and could be removed in future iterations of the system. This would reduce the cost and complexity of the system, and possibly reduce overall friction in the system.

As the results in this thesis are merely based on simulations, they should be interpreted as simulation-based proof of concepts rather than a validation of the final system. Several important effects were simplified or omitted, including nonlinear joint stiffness, how the tendons are routed, cable elasticity, detailed bowden tube friction, exact whiffle-tree geometry, and likely many other factors. Hence, careful implementation and testing on the physical system will be required to validate the results, and to further refine the models and control systems. 